% ---- EDy2: Elektrostatisches Feld ----------------------------

% 2.1 / 2.2  Coulomb & Elektrische Feldstärke
\bereich[cKap2]{2.1 Coulombsches Gesetz \& 2.2 Elektrische Feldstärke $\vec{E}$}
\begin{formelreihe}
  \parbox[t]{\linewidth}{%
    \formelblock{Coulomb-Kraft (Punktladungen)}{%
      \vec{F}=\tfrac{1}{4\pi\varepsilon}\cdot\tfrac{Q_1\cdot Q_2}{r^2}\cdot\vec{e}_r}\\[0.25em]%
    \formelblock{Elektrische Feldstärke}{%
      \vec{E}=\tfrac{\vec{F}}{Q_p}=\tfrac{1}{4\pi\varepsilon}\cdot\tfrac{Q}{r^2}\cdot\vec{e}_r}\\[0.25em]%
    \infoblock{Richtung $\vec{E}$ = Richtung von $\vec{F}$ auf positive Probeladung $Q_p$.\\
      $[\vec{E}]=1\,\tfrac{\text{N}}{\text{As}}=1\,\tfrac{\text{V}}{\text{m}}$}} &
  \parbox[t]{\linewidth}{%
    \formelblock{Verschiebungsdichte (materialunabh.)}{%
      \vec{D}=\varepsilon\cdot\vec{E}\;,\quad
      \left[\vec{D}\right]=\tfrac{\text{As}}{\text{m}^2}}\\[0.25em]%
    \formelblock{Punktladung $Q$}{%
      \vec{D}=\tfrac{Q}{4\pi r^2}\cdot\vec{e}_r}\\[0.25em]%
    \formelblock{Hüllfläche $A$}{%
      Q=\oint_A\!\vec{D}\,\mathrm{d}\vec{A}}\\[0.25em]%
    \formelblock{Divergenz (Gauss)}{%
      \operatorname{div}\vec{D}=\tfrac{\mathrm{d}Q}{\mathrm{d}V}=\rho}} &
  \parbox[t]{\linewidth}{%
    \beispielblock{Feldlinienbild}{%
      \begin{aligned}
        &\text{Richtung}\ \vec{E}=\text{Tangente}\\
        &\text{Betrag}\ |\vec{E}|=\text{Dichte}\\[0.2em]
        &\oint_A\!\vec{E}\,\mathrm{d}\vec{A}\neq0\;\Rightarrow\text{nicht quellfrei}
      \end{aligned}}\\[0.25em]
    \grafik[scale=0.6]{%
      \foreach \k in {0.32,0.72,1.2}{%
        \draw[fld] (-0.7,0) .. controls (-0.28,\k) and (0.28,\k) .. (0.7,0);
        \draw[fld] (-0.7,0) .. controls (-0.28,-\k) and (0.28,-\k) .. (0.7,0);}
      \draw[fld] (-0.55,0)--(0.6,0);
      \node[plus] at (-0.7,0){$+$};
      \node[minus] at (0.7,0){$-$};}} \\[0.3em]
\end{formelreihe}
\bereichende

% 2.3  Spannung und Potential
\bereich[cKap2]{2.3 Spannung und Potential}
\begin{formelreihe}
  \parbox[t]{\linewidth}{%
    \formelblock{Arbeit (homogenes Feld)}{%
      W=Q\cdot E\cdot s\cdot\cos\alpha}\\[0.25em]%
    \formelblock{Homogenfeld: Plattenkond.}{%
      E=\tfrac{U_{12}}{s_{12}}=\tfrac{U}{d}}\\[0.25em]%
    \formelblock{Allgemein (Linienintegral)}{%
      U_{12}=\int_1^2\!\vec{E}\,\mathrm{d}\vec{s}=-U_{21}}} &
  \parbox[t]{\linewidth}{%
    \formelblock{Potential $\varphi$}{%
      \varphi=\tfrac{W}{Q}\;;\quad
      \varphi_a-\varphi_b=U_{ab}=\int_a^b\!\vec{E}\,\mathrm{d}\vec{s}}\\[0.25em]%
    \formelblock{Wirbelfrei (Elektrostatik)}{%
      \oint\vec{E}\,\mathrm{d}\vec{s}=0
      \;\Leftrightarrow\;\operatorname{rot}\vec{E}=0}\\[0.25em]%
    \formelblock{Potentialgradient}{%
      \vec{E}=-\operatorname{grad}\varphi=-\nabla\varphi}} &
  \beispielblock{$\varphi$ im Homogenfeld}{%
    \begin{aligned}
      &\alpha=90°\Rightarrow W=0\\
      &\Rightarrow\text{Äquipotentialflächen}\\
      &\quad\perp\text{ Feldlinien}\\[0.15em]
      &W_{ab}=Q(\varphi_a-\varphi_b)
    \end{aligned}} \\[0.3em]
  \formelblock{Gradient (kartesisch)}{%
    \vec{E}=-\nabla\varphi=-\!\left(\tfrac{\partial\varphi}{\partial x}\vec{e}_x
      +\tfrac{\partial\varphi}{\partial y}\vec{e}_y
      +\tfrac{\partial\varphi}{\partial z}\vec{e}_z\right)} &
  \formelblock{Divergenz $\vec{D}$ (kartesisch)}{%
    \operatorname{div}\vec{D}=\tfrac{\partial D_x}{\partial x}
      +\tfrac{\partial D_y}{\partial y}+\tfrac{\partial D_z}{\partial z}=\nabla\cdot\vec{D}=\rho} &
  \beispielblock{Zusammenfassung Elektrostatik}{%
    \begin{aligned}
      \operatorname{rot}\vec{E}&=0\quad\text{(wirbelfrei)}\\
      \vec{E}&=-\operatorname{grad}\varphi\\
      \operatorname{div}\vec{D}&=\rho\\
      \vec{D}&=\varepsilon\vec{E}
    \end{aligned}} \\[0.3em]
\end{formelreihe}
\bereichende

% 2.4  Kapazität
\bereich[cKap2]{2.4 Kapazität}
\begin{formelreihe}
  \parbox[t]{\linewidth}{%
    \formelblock{Definition}{%
      C=\tfrac{Q}{U}\;,\quad[C]=1\,\tfrac{\text{As}}{\text{V}}=1\,\text{F}}\\[0.25em]%
    \formelblock{Längenbez. Kapazität}{%
      C'=\tfrac{C}{l}=\tfrac{\lambda}{U}\quad\left[\tfrac{\text{As}}{\text{m}^2}\right]}\\[0.25em]%
    \formelblock{Allgemein}{%
      C=\tfrac{Q}{U}=\tfrac{\oint\vec{D}\,\mathrm{d}\vec{A}}{\int\vec{E}\,\mathrm{d}\vec{s}}}} &
  \parbox[t]{\linewidth}{%
    \formelblock{Plattenkondensator}{%
      C=\varepsilon\cdot\tfrac{A}{d}}\\[0.25em]%
    \formelblock{Zylinderkondensator}{%
      \begin{aligned}
        C&=\tfrac{2\pi\varepsilon l}{\ln(r_a/r_i)}\;,\quad
        C'=\tfrac{2\pi\varepsilon}{\ln(r_a/r_i)}\\[0.15em]
        E(r)&=\tfrac{\lambda}{2\pi\varepsilon r}=\tfrac{U}{r\ln(r_a/r_i)}\\[0.15em]
        \varphi(r)&=\tfrac{\lambda}{2\pi\varepsilon}\ln\tfrac{r_a}{r}
      \end{aligned}}} &
  \parbox[t]{\linewidth}{%
    \formelblock{Doppelleitung (Radien $r$, Abstand $a$)}{%
      C=\tfrac{\pi\varepsilon l}{\ln(a/r)}\;,\quad C'=\tfrac{\pi\varepsilon}{\ln(a/r)}}\\[0.25em]%
    \formelblock{Einzelleitung über Erde (Höhe $h$)}{%
      C=\tfrac{2\pi\varepsilon l}{\ln(2h/r)}}\\[0.25em]%
    \infoblock{$C=f(\text{Geometrie, Material})$\\
      $\neq f(U,I,E,D,\ldots)$}\\[0.25em]%
    \grafik[scale=0.72]{%
      % Plattenkondensator
      \draw[line width=0.7pt,red!70!black] (0,0)--(0,0.8);
      \draw[line width=0.7pt,blue!70!black] (0.9,0)--(0.9,0.8);
      \foreach \y in {0.16,0.4,0.64}{\draw[fld] (0.06,\y)--(0.84,\y);}
      \node[font=\tiny] at (0.45,-0.18){Platten};
      % Zylinderkondensator
      \draw[red!70!black] (2.05,0.4) circle(0.42);
      \draw[blue!70!black] (2.05,0.4) circle(0.13);
      \foreach \a in {0,45,...,315}{\draw[fld] ($(2.05,0.4)+(\a:0.15)$)--($(2.05,0.4)+(\a:0.4)$);}
      \node[font=\tiny] at (2.05,-0.18){Zylinder};}} \\[0.3em]
\end{formelreihe}
\bereichende

% 2.5  Energie
\bereich[cKap2]{2.5 Energie im elektrostatischen Feld}
\begin{formelreihe}
  \parbox[t]{\linewidth}{%
    \formelblock{Energiedichte}{%
      w_e=\tfrac{1}{2}\varepsilon E^2=\tfrac{1}{2}D\cdot E=\tfrac{D^2}{2\varepsilon}}\\[0.25em]%
    \formelblock{Energie im Volumen}{%
      W_e=\int w_e\,\mathrm{d}V}} &
  \parbox[t]{\linewidth}{%
    \formelblock{Rotation $\vec{E}$ (Wirbel, kartesisch)}{%
      \operatorname{rot}\vec{E}=\nabla\times\vec{E}=
      \begin{pmatrix}
        \partial_y E_z-\partial_z E_y\\
        \partial_z E_x-\partial_x E_z\\
        \partial_x E_y-\partial_y E_x
      \end{pmatrix}}\\[0.25em]%
    \formelblock{Elektrostatik: wirbelfrei}{%
      W_e=Q\!\oint_{K_A}\!\vec{E}\,\mathrm{d}\vec{s}=0}} &
  \beispielblock{Plattenkond. $A,d$}{%
    \begin{aligned}
      w_e&=\tfrac{1}{2}\varepsilon E^2\\
      W_e&=w_e\cdot A\cdot d\\
         &=\tfrac{1}{2}CU^2
    \end{aligned}} \\[0.3em]
\end{formelreihe}
\bereichende

% 2.6  Grenzflächen
\bereich[cKap2]{2.6 Grenzflächen im elektrostatischen Feld}
\begin{formelreihe}
  \parbox[t]{\linewidth}{%
    \formelblock{Tangentialkomp. stetig (aus $\oint\vec{E}\,\mathrm{d}\vec{s}=0$)}{%
      E_{t1}=E_{t2}}\\[0.25em]%
    \formelblock{Metall ($\kappa\to\infty$)}{%
      E_t=0}} &
  \parbox[t]{\linewidth}{%
    \formelblock{Normalkomp. stetig (aus $\oint\vec{D}\,\mathrm{d}\vec{A}=0$)}{%
      D_{n1}=D_{n2}}\\[0.25em]%
    \formelblock{Brechung der Feldlinien}{%
      \tfrac{\tan\alpha_1}{\tan\alpha_2}=\tfrac{\varepsilon_1}{\varepsilon_2}}} &
  \beispielblock{Grenzfläche $\varepsilon_1|\varepsilon_2$}{%
    \begin{aligned}
      \varepsilon_1 E_{n1}&=\varepsilon_2 E_{n2}\\
      E_{t1}&=E_{t2}\\[0.1em]
      &\Rightarrow\text{Brechung}
    \end{aligned}} \\[0.3em]
\end{formelreihe}
\bereichende

% 2.7  Vektoroperatoren
\bereich[cKap2]{2.7 Vektoroperatoren — Kartesische Koordinaten $(x,y,z)$}
\begin{formelreihe}
  \parbox[t]{\linewidth}{%
    \formelblock{Linien-/Flächen-/Volumenelement}{%
      \begin{aligned}
        \mathrm{d}\vec{s}&=\mathrm{d}x\,\vec{e}_x+\mathrm{d}y\,\vec{e}_y+\mathrm{d}z\,\vec{e}_z\\
        \mathrm{d}V&=\mathrm{d}x\,\mathrm{d}y\,\mathrm{d}z
      \end{aligned}}\\[0.25em]%
    \formelblock{Nabla-Operator}{%
      \nabla=\tfrac{\partial}{\partial x}\vec{e}_x
        +\tfrac{\partial}{\partial y}\vec{e}_y
        +\tfrac{\partial}{\partial z}\vec{e}_z}} &
  \parbox[t]{\linewidth}{%
    \formelblock{Gradient (Skalarfeld $\varphi$)}{%
      \operatorname{grad}\varphi=\nabla\varphi=
        \tfrac{\partial\varphi}{\partial x}\vec{e}_x
        +\tfrac{\partial\varphi}{\partial y}\vec{e}_y
        +\tfrac{\partial\varphi}{\partial z}\vec{e}_z}\\[0.25em]%
    \formelblock{Divergenz}{%
      \operatorname{div}\vec{D}=\nabla\cdot\vec{D}=
        \tfrac{\partial D_x}{\partial x}+\tfrac{\partial D_y}{\partial y}
        +\tfrac{\partial D_z}{\partial z}}\\[0.25em]%
    \formelblock{Laplace}{%
      \Delta=\tfrac{\partial^2}{\partial x^2}
        +\tfrac{\partial^2}{\partial y^2}+\tfrac{\partial^2}{\partial z^2}}} &
  \beispielblock{Bsp: $\varphi=\varphi_0{-}ax$}{%
    \begin{aligned}
      \vec{E}&=-\nabla\varphi=a\,\vec{e}_x\\
      &\Rightarrow\text{homogenes Feld}\\
      &\quad\text{in }x\text{-Richtung}
    \end{aligned}} \\[0.3em]
\end{formelreihe}
\bereichende

\bereich[cKap2]{2.7 Vektoroperatoren — Zylinderkoordinaten $(r,\alpha,z)$}
\begin{formelreihe}
  \parbox[t]{\linewidth}{%
    \formelblock{Linien-/Volumenelement}{%
      \begin{aligned}
        \mathrm{d}\vec{s}&=\mathrm{d}r\,\vec{e}_r+r\,\mathrm{d}\alpha\,\vec{e}_\alpha
          +\mathrm{d}z\,\vec{e}_z\\
        \mathrm{d}V&=r\,\mathrm{d}r\,\mathrm{d}\alpha\,\mathrm{d}z
      \end{aligned}}\\[0.25em]%
    \formelblock{Gradient}{%
      \nabla\varphi=\tfrac{\partial\varphi}{\partial r}\vec{e}_r
        +\tfrac{1}{r}\tfrac{\partial\varphi}{\partial\alpha}\vec{e}_\alpha
        +\tfrac{\partial\varphi}{\partial z}\vec{e}_z}} &
  \parbox[t]{\linewidth}{%
    \formelblock{Divergenz}{%
      \operatorname{div}\vec{D}=
        \tfrac{1}{r}\tfrac{\partial(r D_r)}{\partial r}
        +\tfrac{1}{r}\tfrac{\partial D_\alpha}{\partial\alpha}
        +\tfrac{\partial D_z}{\partial z}}\\[0.25em]%
    \formelblock{Laplace}{%
      \Delta=\tfrac{1}{r}\tfrac{\partial}{\partial r}\!\left(r\tfrac{\partial}{\partial r}\right)
        +\tfrac{1}{r^2}\tfrac{\partial^2}{\partial\alpha^2}
        +\tfrac{\partial^2}{\partial z^2}}} &
  \beispielblock{Zylinderkondensator}{%
    \begin{aligned}
      D(r)&=\tfrac{\lambda}{2\pi r}\\
      E(r)&=\tfrac{\lambda}{2\pi\varepsilon r}\\
      \varphi(r)&=\tfrac{\lambda}{2\pi\varepsilon}\ln\tfrac{r_a}{r}
    \end{aligned}} \\[0.3em]
\end{formelreihe}
\bereichende

\bereich[cKap2]{2.7 Vektoroperatoren — Kugelkoordinaten $(r,\vartheta,\alpha)$}
\begin{formelreihe}
  \parbox[t]{\linewidth}{%
    \formelblock{Volumenelement}{%
      \mathrm{d}V=r^2\sin\vartheta\,\mathrm{d}r\,\mathrm{d}\vartheta\,\mathrm{d}\alpha}\\[0.25em]%
    \formelblock{Gradient}{%
      \nabla\varphi=\tfrac{\partial\varphi}{\partial r}\vec{e}_r
        +\tfrac{1}{r}\tfrac{\partial\varphi}{\partial\vartheta}\vec{e}_\vartheta
        +\tfrac{1}{r\sin\vartheta}\tfrac{\partial\varphi}{\partial\alpha}\vec{e}_\alpha}} &
  \parbox[t]{\linewidth}{%
    \formelblock{Divergenz}{%
      \operatorname{div}\vec{D}=
        \tfrac{1}{r^2}\tfrac{\partial(r^2 D_r)}{\partial r}
        +\tfrac{1}{r\sin\vartheta}\tfrac{\partial(\sin\vartheta\, D_\vartheta)}{\partial\vartheta}
        +\tfrac{1}{r\sin\vartheta}\tfrac{\partial D_\alpha}{\partial\alpha}}\\[0.25em]%
    \formelblock{Laplace}{%
      \Delta=\tfrac{1}{r^2}\tfrac{\partial}{\partial r}\!\left(r^2\tfrac{\partial}{\partial r}\right)
        +\tfrac{1}{r^2\sin\vartheta}\tfrac{\partial}{\partial\vartheta}\!\left(\sin\vartheta
        \tfrac{\partial}{\partial\vartheta}\right)
        +\tfrac{1}{r^2\sin^2\vartheta}\tfrac{\partial^2}{\partial\alpha^2}}} &
  \beispielblock{Punktladung $Q$}{%
    \begin{aligned}
      \vec{D}&=\tfrac{Q}{4\pi r^2}\vec{e}_r\\
      \vec{E}&=\tfrac{Q}{4\pi\varepsilon r^2}\vec{e}_r\\
      \varphi&=\tfrac{Q}{4\pi\varepsilon r}
    \end{aligned}} \\[0.3em]
\end{formelreihe}
\bereichende
